% Author numerical revision; all values derive from actual A/B raw records.
% A separately executed later development phase must retain this evidence.
\section{Original investment decisions: information and complete cost}
\label{sec:study119}

\subsection{The economic comparison and its information budget}
We evaluate whether resolving the active innovation tails reduces the cost of
a certified original Bellman decision.  The economic model remains
\eqref{eq:investment95}--\eqref{eq:cost95}: two investment dates, a nonlinear
transition, the original state-dependent investment capacity, discount
$15/16$, a continuous uniform innovation on $[-1/32,1/32]$, and nested
upper-cost AVaR at $\alpha=1/2$ and $3/4$.  The actual last-date controller,
its observation rule, and its $2^{-20}$ uniform continuation allowance are
unchanged.  At the first date every completed service must certify regret
at most $1/256$ and an original-policy cost interval of width at most
$1/8192$.  Theorem~\ref{thm:investmenttail119} establishes why these
criteria concern the original economic objective, including the implemented
future decisions.  Numerical integration encloses the original continuous
law throughout.

The intervention holds each acquired continuation representation and its
realized first-date proposal fixed while changing the information used for
qualification.  The eight arms are the online semi-analytic Bellman control,
quadratic, ReLU, its equivalent conventional hinge, coordinate spline,
degree-two and degree-three polynomials, and grid interpolation.  The
rank-six quadratic, ReLU and spline constructions retain their common
32-feature dictionaries, three nonlinear selections, nonnegative nonlinear
coefficients, tie rule and 30-bit rounding.  The broader polynomial and
grid baselines retain their additional representational flexibility and
reported memory requirements.  The hinge receives exactly the acquired
ReLU coefficients and its full standalone acquisition charge.  It provides
a conventional implementation of the same information and feasible policies.

In each phase we newly acquire the 256 original-model terminal-value labels
on the fixed $16\times16$ grid for each risk level.  Every competing fit
receives the same label bytes.  Acquisition includes all candidates,
exclusions, independent feature-matrix and active-face QR audits,
serialization and reloading.  The primary catalogue is all 144 states of
the fixed $12\times12$ grid.  These are deterministic numerical catalogues;
no population confidence interval or unseen-sample claim attaches to them.
The experiment concerns the acquired last-continuation slice of NBO in
this original economy.  It does not supply a new trained long-horizon or
empirically calibrated application.

Phase A froze its executable protocol at \texttt{1f33c266} before its new
acquisition and execution.  It compared the complete 128-bin compiled
service (C), the retained geometric whole-support schedule (G), and the
new active-tail schedule (T), producing 6,912 primary decisions and 10,752
timing decisions.  Each service pays for its proposal, every rejected
certificate, refinement and conventional repair.  All 128-bin fallbacks
and eventual exact fallbacks remain available.  The seven fresh timing
processes each use the first 32 catalogue states, with service order
rotated by the declared rule.  All clocks, including the disclosed
document-formatting overlap during phase A, remain archived without
replacement.

Phase A exposed a specific computational problem.  Active-tail qualification
reduced original terminal-risk queries from 654,836 to 293,844 over the
full primary catalogue, but the geometric service required only 223,767.
It beat the compiled service in 12 of 16 median-clock comparisons and the
geometric service in only one; the latter ratio was $0.9947$, with only
three of seven paired clocks lower.  ReLU was slower than both controls
at both risk levels in every recorded repeat.  The full 48-row frontier,
including these losses, remains in the development companion.

\subsection{A bounded preliminary certificate and complete services}
Phase B was designed after observing phase A and prospectively committed
at \texttt{357b2c56} before its scientific execution.  It is disclosed
development on an already observed state catalogue.  It freshly repeated
acquisition, all three controls, reference work and every timing process;
phase-A observations were not used as phase-B controls.  Two additional
services isolate the implementation and precision mechanisms.  The masked
service (M) omits zero-mass tail queries and floating interior searches
after a boundary exit, retaining all six resolutions $2,4,8,16,32,64$.
The budgeted service (B) uses the same masking and original-policy
certificate, but permits only resolutions 2 and 4 before the unchanged
128-bin gate and complete conventional fallback.  It bounds preliminary
outward terminal-value payloads by 30 per decision, or 16 for a boundary
tail, and floating proposal values by 16.  The bounds include the two
outward endpoint payloads that certify the selected tail's true AVaR
defect.  Floating signs never grant certificate authority.

\begin{table}[t]\centering\small
\caption{Phase B: complete 32-state service times, all arms and services}
\label{tab:timing119b}
\begin{tabular}{llrrrrr}\toprule
$\alpha$ & Representation & C & G & T & M & B\\\midrule
1/2 & Online & 0.1015 & 0.0778 & 0.0895 & 0.1107 & 0.1096 \\
1/2 & Quadratic & 0.0588 & 0.0336 & 0.0422 & 0.0650 & 0.0646 \\
1/2 & ReLU & 0.0896 & 0.0965 & 0.1187 & 0.1361 & 0.1071 \\
1/2 & Hinge & 0.0819 & 0.0865 & 0.1097 & 0.1256 & 0.0957 \\
1/2 & Spline & 0.0582 & 0.0318 & 0.0417 & 0.0611 & 0.0604 \\
1/2 & Polynomial 2 & 0.2227 & 0.2022 & 0.2170 & 0.2353 & 0.2298 \\
1/2 & Polynomial 3 & 0.2864 & 0.2509 & 0.2665 & 0.2901 & 0.2858 \\
1/2 & Grid & 0.2343 & 0.1983 & 0.2091 & 0.2342 & 0.2269 \\
3/4 & Online & 0.1030 & 0.0761 & 0.0761 & 0.0881 & 0.0873 \\
3/4 & Quadratic & 0.0584 & 0.0318 & 0.0324 & 0.0431 & 0.0429 \\
3/4 & ReLU & 0.0904 & 0.0969 & 0.1182 & 0.1189 & 0.0934 \\
3/4 & Hinge & 0.0801 & 0.0848 & 0.1076 & 0.1119 & 0.0835 \\
3/4 & Spline & 0.0585 & 0.0317 & 0.0324 & 0.0419 & 0.0430 \\
3/4 & Polynomial 2 & 0.2190 & 0.1853 & 0.1870 & 0.1984 & 0.2182 \\
3/4 & Polynomial 3 & 0.2655 & 0.2398 & 0.2420 & 0.2515 & 0.2485 \\
3/4 & Grid & 0.2262 & 0.1936 & 0.1981 & 0.2071 & 0.2127 \\
\bottomrule\end{tabular}
\par\smallskip\footnotesize
Medians in seconds from all seven fresh sequential processes. C: 128-bin compiled; G: geometric whole-support; T: original active tail; M: positive-tail masking with the full ladder; B: masked two-level budget. Every service pays its proposal, refusals, repairs and certificates. Acquisition, actual loading and enclosing process costs are reported separately. Observed ranges and every clock are retained in the development companion.
\end{table}


All 11,520 primary and 17,920 timing decisions completed.  Independent
record arithmetic verified 1,031,933 claims over 45,557 attempted
certificates.  A separate read-only reconstruction checked all 640 raw
file hashes and 295,161 additional qualification, policy-identity and
payload-account claims.  The two-level budget attained but never exceeded
30 outward values, and attained but never exceeded 16 at a boundary,
over all 5,888 primary and timing rows of that service.  It used 28,801
preliminary outward payloads in the primary catalogue, compared with
94,925 for the full masked ladder.  The 244 primary proposals exhausting
that ladder were the same 244 requiring the complete fallback under the
two-level rule.  All were retained and paid.

Table~\ref{tab:timing119b} reports every arm, risk and service.  Masking
alone reduced risk queries from 293,844 to 220,341, yet was slower than
the original active-tail implementation in all 16 median comparisons.
The budgeted rule reduced them further to 154,217, but beat geometric
whole-support qualification in only two of 16 medians.  These were ReLU
and the equivalent hinge at $\alpha=3/4$, with ratios $0.9640$ and
$0.9843$; their seven paired-clock win counts were seven and five.
ReLU at $\alpha=1/2$ instead had ratio $1.1093$ and lost in all seven
repeats.  Thus lower information demand and fewer preliminary resolutions
do not establish a generally cheaper implementation.  Mask grouping,
interval assembly, recorded traces, proposal inference and complete
fallback remain actual service work.

The 144-state workload gives a related but different comparison.  Budgeted
qualification was faster than geometric qualification in seven of 16
cells and faster than compiled qualification in all 16.  Summing the
observed primary service times across all arms and risks gives $9.2482$,
$7.2390$, $8.0358$, $8.6385$, and $7.5186$ seconds for C, G, T, M, and B.
These totals retain all cells and are not a selectively composed service.
The small-batch and full-catalogue rankings need not agree: the first 32
states contain a different incidence of costly initial proposals, and
vectorized grouping has a workload-dependent cost.  We do not pool their
clocks or assign an unobserved request distribution.

\subsection{Original policy costs, action error and adoption}
The reference calculation partitions the full original innovation support
into 1,024 bins independently of the qualification mesh and evaluates every
distinct returned state-action policy. It shares the verified terminal
compiler with production; it is not a fully independent implementation.
The separately implemented exact-rational primitive suite reconstructs
24 original reference cases and 72 certificates, and has that finite scope.  It pays for 983 distinct policy queries and 9.6771
seconds in phase B.  The best valid all-action lower bound and feasible
policy upper bound at each state enclose the original optimum.  They
give original-policy regret intervals and action-distance bounds using
original action curvature two.  The largest reported regret upper bound
is $0.00285149$, below $1/256$; the largest action-distance upper bound is
$0.053400$.  These are verified bounds, not exact measured action errors.

Every one of the ten within-arm service comparisons has exactly the same
returned complete policy at all 2,304 state--risk--arm combinations.  Its
original economic cost contrast is therefore exactly zero.  The 1,220
recorded primary action repairs are 244 distinct initial proposals paid
and repaired separately under five services.  No exact rational fallback
was required.  More informative certificates change the information
required to serve these policies; this experiment supplies no new policy
improvement event.  The distinct original-policy improvements of the
retained R110 safeguard and gain-budget campaigns remain reported with
their own sources, catalogues and clocks in Section~\ref{sec:study110}.

\begin{table}[t]\centering\small
\caption{Phase B: paid acquisition, complete 144-state cost and economic bounds}
\label{tab:cost119b}
\begin{tabular}{llrrrrr}\toprule
$\alpha$ & Representation & $A$ & $F_G$ & $F_B$ & $10^6\overline L$ & $10^3\overline d_{\max}$\\\midrule
1/2 & Online & 0.0001 & 0.2960 & 0.3496 & 0.899 & 0.957 \\
1/2 & Quadratic & 0.0634 & 0.1648 & 0.1874 & 153.657 & 31.496 \\
1/2 & ReLU & 0.0599 & 0.3927 & 0.3755 & 171.989 & 51.389 \\
1/2 & Hinge & 0.0602 & 0.3341 & 0.3182 & 171.989 & 51.389 \\
1/2 & Spline & 0.0536 & 0.3274 & 0.3395 & 264.595 & 53.400 \\
1/2 & Polynomial 2 & 0.0105 & 0.6376 & 0.8062 & 153.657 & 31.495 \\
1/2 & Polynomial 3 & 0.0107 & 0.7882 & 0.9014 & 215.969 & 29.648 \\
1/2 & Grid & 0.0101 & 0.8293 & 0.7160 & 11.069 & 11.106 \\
3/4 & Online & 0.0002 & 0.2823 & 0.3395 & 0.900 & 1.022 \\
3/4 & Quadratic & 0.0698 & 0.1620 & 0.1796 & 163.145 & 31.487 \\
3/4 & ReLU & 0.0638 & 0.4118 & 0.3642 & 181.578 & 49.019 \\
3/4 & Hinge & 0.0639 & 0.3765 & 0.3436 & 181.578 & 49.019 \\
3/4 & Spline & 0.0603 & 0.3851 & 0.3464 & 289.860 & 52.111 \\
3/4 & Polynomial 2 & 0.0154 & 0.7420 & 0.8229 & 163.147 & 31.487 \\
3/4 & Polynomial 3 & 0.0155 & 0.8937 & 0.9413 & 209.222 & 30.155 \\
3/4 & Grid & 0.0150 & 0.7887 & 0.7603 & 35.811 & 18.471 \\
\bottomrule\end{tabular}
\par\smallskip\footnotesize
$A$ includes original labels, every fit and candidate audit, and serialization; hinge includes the full ReLU acquisition. $F_s=A+\text{actual load time}+\text{actual 144-state service time}$, in seconds. These are recorded component sums, not new clocks. $\overline L$ is an outward upper bound for catalogue-mean original-policy regret; $\overline d_{\max}$ bounds the largest distance to the original optimal action. All five services have the same returned policies, so their economic bounds coincide. Displayed bound endpoints are rounded outward.
\end{table}


The conventional frontier is economically material.  ReLU's mean original
cost exceeds quadratic's by an interval of
$[1.79075,1.87557]\times10^{-5}$ at $1/2$ and
$[1.80271,1.88385]\times10^{-5}$ at $3/4$.
ReLU is individually cheaper at 43 and 39 states, but the whole-catalogue
mean differences are strictly adverse.  Against online control, it is
identical at 108 states and strictly costlier at 36 under either risk
level.  Its mean cost exceeds grid interpolation's and is below spline's
and degree-three polynomial's; every pairwise interval, including
unresolved statewise comparisons, is retained.  The exact ReLU--hinge
identity holds for 80 primary/timing record pairs containing 3,680 paired
decisions and all their attempted certificates.  Their differing elapsed
times cannot be attributed to an exclusive neural representation.

Table~\ref{tab:cost119b} includes each arm's actual standalone acquisition
and complete 144-state cost.  Quadratic with geometric qualification costs
$0.1648$ and $0.1620$ seconds, versus $0.3755$ and $0.3642$ for the
budgeted ReLU service.  It also has the lower certified catalogue-mean
economic cost.  For this entire 144-state workload the ReLU acquisition
saving is insufficient to offset service work.  This observation leaves
neural adoption unestablished on the measured frontier.  The inherited
least-squares acquisition criterion also remains economically imperfect:
the degree-three fit has lower terminal fitting loss than quadratic,
while its original catalogue-mean regret upper bound is higher.  New
certificates do not turn this into action-aligned training.

For a deployment repeating this same 144-state catalogue, let $A_j$ include
acquisition and one model load, $C_j$ the complete per-catalogue service
cost, and $\Delta J$ the mean original-policy cost difference from a
specified comparator.  With $H$ catalogue repetitions and a specified
value $\lambda\ge0$ per economic-cost unit in the same accounting units,
a sufficient adoption condition is
\[
 A_j-A_c+H(C_j-C_c)+144H\lambda\,\overline{\Delta J}<0.
\]
It uses an upper economic contrast bound, not a difference of regret
ceilings.  No $H$, compute price, monetary calibration or request
population is supplied by the experiment.  At fixed identical policies
the economic term is zero.  Measured information savings therefore
justify adoption only when complete implementation cost, acquisition and
the intended workload satisfy the corresponding condition.

\subsection{Information theory and the scope of scaling evidence}
The separate exact point-information catalogue tests
Theorem~\ref{thm:tailinformation119} on ten predeclared convex
piecewise-linear/quadratic functions, eight risk levels from zero through
$255/256$, and four resolutions.  All 320 cells and 2,901 rational checks
are retained.  The active-tail interval is narrower in 242 cells,
whole-support quadrature is narrower in 27, and 51 tie.  Both methods
obey the same declared query cap; actual distinct queries are reported
separately.  At $\alpha=0$, whole-support wins 20 of 40 cells and active
integration wins eight.  Nearer the upper risk limit active localization
uses the shorter adverse support, consistently with the theorem's
$1-\alpha$ factor.  Exact independent integrals and cut reconstruction
are evaluation evidence; formula access is not provided to either
point-query algorithm.

Phase A also compares the same original online proposal with a compiled
128-bin certificate and an eight-cell-per-tail certificate at dimensions
$2,4,8$.  All 192 paired contexts qualify.  The active/compiled certificate
time ratios are approximately $0.515,0.475$, $0.423,0.424$, and
$0.391,0.317$ across the two risks.  Charging their common proposal moves
the ratios to $0.829,0.793$, $0.757,0.763$, and $0.690,0.659$.
These are original-model primitive scaling observations.  They do not
establish a complete learned high-dimensional service, a multivariate
risk integration theorem, or a longer-horizon convergence rate.  The
mathematical information improvement, its action-valid certificate, and
these implementation limitations together specify the next economic
computational problem; all historical adverse and incomplete episodes
remain attributable in the development companion.
